\subsection{Homomorphic Arithmetic Aggregation}\label{block: LWEsPack}
To support homomorphic arithmetic aggregations 
(such as $\dbsum$, $\dbcount$, and $\dbavg$), 
the $\dbtabrowsize$ LWE ciphertexts which encrypt the filtered results
have to be packed in an RLWE ciphertext 
as arithmetic aggregation over LWE ciphertexts are slow.
In this section, we propose a new homomorphic ciphertext conversion
algorithm $\repack$ and the associated 
homomorphic arithmetic aggregations procedure.

\begin{algorithm}[t!]
    % \tiny
    \caption{$\repack$}
    \label{algo: lwespack}
    \SetKwInOut{Input}{Input}\SetKwInOut{Output}{Output}\SetKwInOut{Return}{Return}\SetKwComment{Comment}{$\triangleright$\ }{}
    \Input{$N$ LWE ciphertexts $(\lwe_0, \lwe_1, ...,\lwe_{N-1})$ where 
    $\lwe_i = \lwe^{n,q}_{\mathbf{s}}(m_i) = (b_i, \mathbf{a_i})$
    and $m_i = 0/1$. }
    \Input{A conversion key 
    $\reg{\mathbf{CK}}$.}
    \Input{A precision parameter $\delta$, range parameter $(u,v)$}
    \Output{An RLWE ciphertext $ct = \rlwe^{N,Q}_{\Tilde{s}}(\encode(\mathbf{m}))$,
    where $\mathbf{m}_{i} = m_i$ for $i \in \mathbb{Z}_{N}$}
     Let $\mathbf{b} \leftarrow [b_0, b_1, ...,b_{N-1}]$, 
     $\mathbf{A} \leftarrow [\mathbf{a_0}, \mathbf{a_1}, ..., \mathbf{a_{N-1}}]^{T} \in \mathbb{Z}_{q}^{N \times n}$\\
     % $ct_1 = \texttt{HomMVMul}(\mathbf{A},CK)$.
     $ct_1 \leftarrow \mathbf{CK} \diamond \mathbf{A}$
     \Comment*[r]{$ct_1 = \rlwe^{N,Q_1}_{\Tilde{s}}(\encode(\mathbf{A}\mathbf{s}))$}
     $ct_2 \leftarrow ct_1 + (\encode(\mathbf{b},0))$
     \Comment*[r]{$ct_2 = \rlwe^{N,Q_1}_{\Tilde{s}}(\encode(\mathbf{A}\mathbf{s} + \mathbf{b}))$} 
     $ct_3 \leftarrow\texttt{HomMod}(ct_2, q)$
     \Comment*[r]{$ct_3 = \rlwe^{N,Q_2}_{\Tilde{s}}(\encode(\mathbf{m+e}))$} 
     Generate the approximate polynomial $p(x)$ \\
     $ct \leftarrow \prnd(ct_3)$
     \Comment*[r]{$ct = \rlwe^{N,Q_3}_{\Tilde{s}}(\encode(\prnd(\mathbf{m})))$} 
    \Return{$ct = \rlwe^{N,Q}_{\Tilde{s}}(\encode(\prnd(\mathbf{m})))$}
\end{algorithm}

\subsubsection{Homomorphic ciphertext conversion $\repack$}
Existing works in converting LWE ciphertexts into
an RLWE ciphertext
include PEGASUS~\cite{pegasus} and 
the method of Chen \etal~\cite{conversionhaochen}.
% Both their methods can only achieve one competitive performance 
% in latency and precision.
% The key idea of 
% Pegasus~\cite{pegasus} in ciphertext conversion 
% consists of homomorphically evaluating LWE partial decryption as well as 
% homomorphically evaluating modulo $q$ on an RLWE ciphertext encrypting secret key.
% Both two operations can be achieved 
% by arithmetic homomorphic circuits~\cite{helib,ckksbootsparse,ckksbootvariance}.
Taking advantage of SIMD packing arithmetic homomorphic circuits, 
PEGASUS~\cite{pegasus}
performs fast ciphertext conversion but suffers from the low-precision 
problem due to the noise growth in homomorphic arithmetic circuits.
In contrast, Chen \etal~\cite{conversionhaochen} packs the
LWE ciphertexts by evaluating a series of
recursive homomorphic automorphisms~\cite{helib}.
We point out that Chen \etal~\cite{conversionhaochen}
enjoys high-precision packing capability, but 
induces a higher amount of latency overheads.
% This approach achieves high-precision packing but is 
% limited to a small size.
% When the packing size of LWE ciphertexts is large, 
% the latency is still very high.

%Unfortunately, in an encrypted database setting, 
Here, we design a homomorphic $\repack$ algorithm that
achieves both, a fast algorithm that can accurately
pack a set of LWE ciphertexts into one RLWE ciphertext.
%and high-precision format conversion.
% In \NAME, we propose our homomorphic ciphertext conversion method
% for homomorphically encrypted database
% which achieves low latency and high precision.
Our critical insight is that, in a homomorphically encrypted database, 
the predicate results are in a very limited range, e.g., usually only
0's and 1's to express the logical validity of the database entries.
Thus, we can actually ``correct'' the errors on the low-precision ciphertext to 
obtain a high-precision ciphertext. 
% To be more specific, we first
% execute the fast repack algorithm proposed in~\cite{pegasus} on
% a set of LWE ciphertexts to produce

% Then, we evaluate an additional $\Round$ function on the repacked RLWE
% ciphertext to homomorphically remove the noise $e$. 
% % the objective of our new conversion algorithm is to convert 
% %$ct = \rlwe(\encode(\mathbf{m}))\text{ where } $\bm \ni m_i = 0\pm e\text{ or }1\pm e$
% Concretely, we wish to convert
% $ct=\rlwe(\encode(\mathbf{m}))$ into $ct' =  \rlwe(\encode(\mathbf{m'}))$ such that 
% $m'_i = \Round(m_i)$, and that
% $m'_i$ contains less noises than $m_i$ for all $i\in|\bm|$.


The detailed algorithm $\repack$ involves the following steps.
First of all, in \cref{algo: lwespack},
we are given $N$ LWE filtering results $\lwe_0, \lwe_1, ...,\lwe_{N-1}$
where $\lwe_i = \reg{LWE}_{\mathbf{s}}^{n,q}(m_i) = (b_i, \mathbf{a_i})$
with a  conversion key 
$\reg{\mathbf{CK}}$.
On Line 1, we follow~\cite{pegasus} and rearrange the ciphertexts to 
construct a ciphertext matrix
$\mathbf{A} = [\mathbf{a_0}, \mathbf{a_1}, ..., \mathbf{a_{N-1}}]^{T} 
\in \mathbb{Z}_{q}^{N \times n}$ and a ciphertext vector 
$\mathbf{b} = [b_0, b_1, ...,b_{N-1}] \in \mathbb{Z}_{q}^{N}$.
Next, on Line 2, we evaluate the 
homomorphic matrix multiplication
$\mathbf{CK} \diamond \mathbf{A}$ 
to get $ct_1 = \rlwe^{N,Q_1}_{\Tilde{s}}(\encode(\mathbf{A}\mathbf{s}))$,
where $\encode$ stands for the CKKS encoding 
function~\cite{DBLP:conf/asiacrypt/CheonKKS17}.
Then, on Line 3, we homomorphically add $\bb$ to $ct_1$ and get $ct_2 = \rlwe^{N,Q_1}_{\Tilde{s}}(\encode(\mathbf{A}\mathbf{s} + \mathbf{b}))$.
After the above transforms, we obtain the ciphertext $ct_2$
encrypting $\mathbf{A}\mathbf{s}+\mathbf{b} = \mathbf{m} + t\mathbf{q}$.
To remove the $t\mathbf{q}$ term, on Line 4, 
we follow~\cite{ckksbootstrap,ckksbootsparse, ckksbootvariance, sinemod} 
to apply \texttt{HomMod} to the \texttt{mod} $q$ function 
homomorphically on $ct_2$, and get
$ct_3$. %i$ = \rlwe^{N,Q_2}_{\Tilde{s}}(\encode(\mathbf{m+e}))$.
At this stage, we obtain a packed but imprecise $\rlwe$ ciphertext 
$ct_3 = \rlwe(\encode(\mathbf{m}))$ encrypting
the vector of messages $\bm$ packed from the input. Here, we observe that for each $m_i\in\bm$,
$m_i$ equals either $1\pm e$ or $0\pm e$, where the noise $e$ 
in inevitably added via the above transforms.

To correct the errors, on Line 5 in \cref{algo: lwespack}, 
we generate a polynomial $\prnd(x)$ to approximate the rounding function $\Round(x)$.
% Similar to the bootstrapping function, to evaluate an approximate
% homomorphic rounding function over RLWE ciphertexts, we need
% to construct a polynomial that approximates the behavior
% of a rounding function.
%As it is not easy to  evaluate a round function
%via homomorphic arithmetic circuits directly, 
%we generate a high-degree polynomial to 
% approximate the target round function.
Following~\cite{ckksbootsparse}, given an interval $[u,v]$ and a pre-defined 
precision parameter $\delta$, we want to obtain
a polynomial $\prnd(x) =  {\textstyle \sum_{i=0}^{n-1}}a_ix^i $ such that
\begin{equation}
\forall x \in [u,v],  |\prnd(x)-\round{x} | < 2^{-\delta}, \label{formula: appro}
\end{equation}
where $2^{-\delta}$ is negligible.
%This question is a classical problem in approximation theory.
Based on the equioscillation theorem~\cite{equioscillation}, $\prnd(x)$
is the best possible approximation of the round function
if and only if 
\begin{align}
&\exists x_0 \leq ... x_i ... \leq x_n \in [u,v], &\\
&\prnd(x_i) - \round{x}= \epsilon(-1)^{i}|\prnd(x) - \round{x}|_{\infty} &
\label{formula: equioscillation}
\end{align}
where $\epsilon$ = $\pm 1$.
Taking \cref{formula: equioscillation} as the judgment condition, 
we can construct the Remez iteration algorithm~\cite{remez} to 
obtain the polynomial $p(x)$ meets the requirement of \cref{formula: appro}.
Upon obtaining $\prnd(x)$, we can actually perform
the approximate homomorphic rounding function on Line 6. As
explained in \cref{pre:he_circuits}, $\prnd(x)$ can be directly applied
to the RLWE $ct_3$ produced on Line 4, and we get the final output
$ct =  \rlwe(\encode(\prnd(\bm)))$. Here, the guarantee
is that the noises in $\prnd(\bm)$ will strictly be smaller than 
that in $\bm$.

% Concretely, we wish to convert
% $ct=\rlwe(\encode(\mathbf{m}))$ into $ct' =  \rlwe(\encode(\mathbf{m'}))$ such that 
% $m'_i = \Round(m_i)$, and that
% $m'_i$ contains less noises than $m_i$ for all $i\in|\bm|$.
% as claimed.
% After generating the approximated round function $p(x)$, we can
% evaluate $p(x)$ homomorphically over the RLWE ciphertext $ct_3$,
% as explained in \cref{pre:he_circuits}.
% As $p(x) =  {\textstyle \sum_{i=0}^{n-1}}a_ix^i $ is a type of 
% polynomial functions, it could be implemented efficiently 
% by homomorphic arithmetic circuits as we introduced in \cref{pre:he_circuits}.
% Moreover, expressing the approximated function $p(x)$ 
% in a Chebyshev basis~\cite{chebyshevbasis}
% could optimize the evaluation process with a less 
% number of ciphertext multiplication~\cite{
% improvedboot, betterboot, ckksbootsparse, inverseboot}.  

\subsubsection{Homomorphic arithmetic aggregation}
Using the filtered results packed in an RLWE ciphertext $cf$ based on
the proposed $\repack$ operator described above, we can finally
aggregate over the target attribute.  Let the RLWE ciphertext
encrypting items in the column be $ct_a$, we describe how
to implement three types of conventional arithmetic aggregation functions, 
namely, $\dbcount$, $\dbsum$ and $\dbavg$.
%encrypt the vector of filtered result and the encrypted column of 
%aggregated attribute $ct_a$.
Note that the first two functions can be implemented using simple inner
products between ciphertexts. Concretely,
$\dbcount$ function can be implemented by the inner product between $cf$ and
the vector $I = (1, 1, ...,1)$, while $\dbsum$
is also an inner product between $cf$ and $ct_a$.
The homomorphic inner product of two RLWE ciphertexts can
be implemented over ciphertexts in both the slot format~\cite{gazelle, pegasus} 
or in the coefficient format~\cite{cheetah, heda}.
Furthermore, the $\dbavg$ function can be evaluated by 
%$\dbavg() = \frac{\dbsum()}{\dbtabrowsize} \times 
% (\frac{\dbcount)}{\dbtabrowsize})^{-1}$, in which
$\dbavg() = \dbsum()\times 
\dbcount()^{-1}$, where the inverse function 
%the $\frac{1}{\dbtabrowsize}$ is public and 
%the inversion function $(\frac{\dbcount()}{\dbtabrowsize})^{-1}$ 
can be implemented based on ~\cite{ckkscomparison}.

% The homomorphic modular reduction function is a central part
% of CKKS bootstrapping~\cite{ckksbootstrap}, 
% %which is a cryptography puzzle of the FHE research community.
% and a number of recent works~\cite{ckksbootsparse, ckksbootvariance, sinemod}
% have investigated the accuracy-efficiency trade-off for different
% approximation polynomials.
% Note that the exact
% homomorphic matrix multiplication algorithms can be found in~\cite{gazelle, pegasus}.
% To add the vector $\mathbf{b}$, 
% we add the encoding result $\encode(\mathbf{b})$ of the vector 
% into the first term 
% of the ciphertext $ct_1$
% and gets $ct_2 = \rlwe^{N,Q_1}_{\Tilde{s}}(\encode(\mathbf{A}\mathbf{s} + \mathbf{b}))$.
% encode the vector $\mathbf{b}$ into
% a polynomial $\encode(\mathbf{b})$ and add it into the first term 
% of the ciphertext $ct_1$.
% Followed by a homomorphic addition with the vector $\mathbf{b}$, 
% we perform the LWE decryption formula $\mathbf{A}\mathbf{s}+\mathbf{b}$
% except for the modulo function.
% While there is a vast amount of prior work~\cite{helib, gazelle, haochenONN, pegasus}
% on HE for matrix-vector multiplication, 
% none take a generally efficient approach in all applications.
% Need to discuss

%\step{Line 4, Modulo $q$}:

%propose their own algorithm to solve it.
% and we adapt the approach from ~\cite{ckksbootsparse}
% as this is not the research part of this paper.